\subsection{Explicación del comando \texttt{mpirun} y \texttt{mpiexec}}

El comando \texttt{mpirun} constituye la herramienta fundamental empleada en entornos de computación de alto rendimiento para iniciar y gestionar aplicaciones distribuidas basadas en MPI. A diferencia de los programas secuenciales convencionales, las aplicaciones paralelas requieren un gestor de procesos coordinado que inicialice múltiples instancias del mismo binario a través de una arquitectura de hardware específica \cite{gropp_using_mpi_1999}. Este comando actúa como una interfaz de lanzamiento universal que abstrae la complejidad de la comunicación inicial entre los nodos de cómputo y el sistema operativo local \cite{pacheco_2011}.

Aunque históricamente la comunidad utilizaba diversas variantes de \texttt{mpirun}, el estándar MPI introdujo formalmente \texttt{mpiexec} para proveer un mecanismo uniforme y portable que evite la confusión derivada de opciones no estandarizadas \cite{mpi_forum_2025}. En la práctica moderna, ambos comandos suelen coexistir como sinónimos funcionales en la mayoría de las implementaciones, aunque \texttt{mpiexec} se prefiere cuando se busca el máximo apego a la especificación oficial. El propósito primordial de este mecanismo de lanzamiento consiste en establecer el entorno de ejecución necesario para que los procesos puedan identificarse y comunicarse entre sí desde el momento de su inicialización, asignando a cada uno un identificador único conocido como rango dentro del comunicador global \texttt{MPI\_COMM\_WORLD} \cite{mpi_forum_2025, grama_intro_parallel_2003}.

La invocación básica requiere la especificación explícita de la cantidad de tareas y la ruta del binario previamente compilado. En su forma más elemental, la sintaxis consiste en definir el parámetro de conteo de procesos seguido inmediatamente por el nombre del ejecutable, como en el comando \texttt{mpirun -np 4 ./programa\_paralelo}, el cual instruye al gestor de la biblioteca a instanciar exactamente cuatro copias independientes del software en los núcleos disponibles \cite{pacheco_2011}. Con el fin de flexibilizar la ejecución en arquitecturas distribuidas complejas, existen múltiples parámetros opcionales orientados a controlar la distribución y el mapeo detallado de la carga de trabajo, los cuales se resumen en el Cuadro~\ref{tab:parametros_mpirun}.

\begin{table}[h]
\centering
\caption{Parámetros opcionales recurrentes en los comandos \texttt{mpirun} y \texttt{mpiexec}.}
\label{tab:parametros_mpirun}
\begin{tabular}{|l|p{10cm}|}
\hline
\textbf{Parámetro} & \textbf{Descripción y propósito} \\ \hline
\texttt{-np <N>} / \texttt{-n <N>} & Especifica la cantidad total de procesos o instancias que se crearán para la aplicación distribuida. \\ \hline
\texttt{-hostfile <archivo>} & Define la ruta hacia un archivo de texto que contiene la lista de nodos y las ranuras de cómputo disponibles. \\ \hline
\texttt{-host <nombres>} & Indica los nodos o máquinas de la red donde específicamente se deben ejecutar los procesos. \\ \hline
\texttt{-arch <arquitectura>} & Permite solicitar nodos que posean una arquitectura de hardware específica (por ejemplo, \texttt{x86\_64}). \\ \hline
\texttt{-wdir <directorio>} & Establece el directorio de trabajo para los procesos instanciados. \\ \hline
\texttt{-configfile <archivo>} & Permite proporcionar un archivo de texto con configuraciones detalladas cuando los argumentos de línea de comandos son demasiados o complejos. \\ \hline
\texttt{-map-by <estrategia>} & Modifica la estrategia de asignación de procesos, permitiendo agruparlos por nodo, ranura o socket. \\ \hline
\texttt{-x <var>} / \texttt{--env <var>} & Exporta variables de entorno específicas del sistema local hacia los entornos de los nodos remotos. \\ \hline
\end{tabular}
\end{table}

El parámetro estructural \texttt{-hostfile} adquiere una relevancia crítica cuando se transiciona de un entorno de pruebas mononodo a un clúster de computadoras interconectadas por red, ya que distribuye de manera balanceada o secuencial las instancias del binario entre las máquinas físicas especificadas \cite{grama_intro_parallel_2003}. Por otro lado, modificadores avanzados como \texttt{-map-by} y \texttt{-x} garantizan el control granular sobre la afinidad del procesador y la propagación de contextos de ejecución homogéneos en sistemas heterogéneos \cite{open_mpi_2026}. Adicionalmente, el comando \texttt{mpiexec} permite la ejecución de aplicaciones heterogéneas bajo el modelo MPMD (\textit{Multiple Program, Multiple Data}), lo cual significa que se pueden iniciar simultáneamente varios programas distintos, o el mismo programa con diferentes argumentos, separando las especificaciones con el símbolo de dos puntos (\texttt{:}) en la línea de comandos \cite{mpi_forum_2025}.

Un aspecto didáctico de alta relevancia para quienes inician en MPI es el comportamiento no determinista de la salida estándar (\texttt{stdout}). Es erróneo asumir que las sentencias de impresión aparecerán en la pantalla ordenadas según el rango de cada proceso, pues los múltiples procesos compiten de forma asíncrona por el acceso al dispositivo de salida, generando una salida desordenada o interrumpida mutuamente \cite{pacheco_2011}. Para evitar esta confusión, se debe estructurar el código de manera que un solo proceso (por ejemplo, el rango 0) recolecte la información de los demás y se encargue exclusivamente de imprimirla en pantalla.

La implementación exitosa de entornos paralelos mediante \texttt{mpirun} o \texttt{mpiexec} demanda una supervisión constante sobre la compatibilidad de las versiones de software y la asignación realista de recursos físicos. Un error recurrente consiste en intentar ejecutar un número de procesos significativamente mayor al número de núcleos disponibles, práctica conocida como sobreasignación o \textit{oversubscription}, que provoca una degradación severa en el rendimiento debido a la latencia inducida por el cambio de contexto excesivo \cite{gropp_using_mpi_1999}. Asimismo, resulta imperativo vigilar la correspondencia exacta entre la implementación utilizada para compilar el código y la herramienta empleada para su ejecución; si se genera un binario con \texttt{mpicc} de MPICH e intentar ejecutarlo con \texttt{mpirun} de Open MPI, se desencadenarán fallos críticos de enlace en tiempo de ejecución o comportamientos indeterminados \cite{open_mpi_2026}. Por lo tanto, la validación previa de los módulos activos y el apego a la sintaxis estándar de \texttt{mpiexec} constituyen las mejores prácticas para mitigar fallas en la computación distribuida y preservar la portabilidad de los \textit{scripts} de lanzamiento \cite{pacheco_2011, mpi_forum_2025}.
